% ============================================================
%  Section 7: Discussion and Impact
% ============================================================

\section{Discussion and Impact}
\label{sec:discussion}

\subsection{Key Contributions}
\label{subsec:contributions}

\bert{} introduced three innovations that, taken together, established a new paradigm
for NLP research and engineering:

\begin{enumerate}
  \item \textbf{Deep bidirectional pre-training.}  By employing the MLM objective,
        \bert{} learns representations that incorporate context from both directions
        simultaneously in every layer.  This is a departure from the shallow
        bidirectionality of ELMo (which concatenates unidirectional LMs) and from
        the strict left-to-right conditioning of GPT.

  \item \textbf{Unified pre-train–fine-tune framework.}  A single model can be
        adapted to diverse tasks --- classification, question answering, sequence
        labelling, semantic similarity --- with minimal architectural change,
        reducing the need for highly specialised task-specific designs.

  \item \textbf{State-of-the-art empirical results.}  Upon publication, \bert{}
        achieved new records on 11 NLP benchmarks, including a 7.7-point
        improvement on the GLUE benchmark~\cite{wang2018glue} and a 1.5 F1-point
        gain on SQuAD~v1.1~\cite{rajpurkar2016squad}.
\end{enumerate}

\subsection{Advantages}
\label{subsec:advantages}

\begin{itemize}
  \item \textbf{Strong transfer learning.}  Pre-trained \bert{} representations
        encode broad linguistic knowledge that transfers well to low-resource tasks
        where labelled data is scarce.

  \item \textbf{Task flexibility.}  The same encoder is applicable to sentence-level
        and token-level tasks without structural modification.

  \item \textbf{Contextual disambiguation.}  Because each token's representation
        is a function of all other tokens in the sequence, polysemous words are
        disambiguated naturally --- a capability absent from static embeddings.
\end{itemize}

\subsection{Limitations}
\label{subsec:limitations}

Despite its impact, \bert{} carries several notable limitations:

\begin{itemize}
  \item \textbf{Computational cost.}  Pre-training \bert{}-Large requires thousands
        of TPU hours and is prohibitively expensive for most academic and industrial
        teams without specialised hardware.  Inference latency is also high relative
        to simpler models, hampering deployment in latency-sensitive applications.

  \item \textbf{Fixed input length.}  The 512-token limit imposed by learned
        position embeddings makes \bert{} unsuitable for long documents without
        chunking heuristics.

  \item \textbf{NSP objective limitations.}  The NSP task does not require
        fine-grained linguistic understanding and can be learned by superficial
        topic-matching heuristics~\cite{liu2019roberta}.

  \item \textbf{Masked token proportion.}  Only 15\% of tokens are masked per
        step, yielding a sparse training signal relative to autoregressive models
        that compute loss over all tokens.

  \item \textbf{Discrepancy between pre-training and fine-tuning.}  The \mask{}
        token appears during pre-training but never during fine-tuning, creating a
        distribution mismatch that may degrade performance.
\end{itemize}

\subsection{Influence and Subsequent Work}
\label{subsec:subsequent}

\bert{}'s design has spawned a rich ecosystem of variants and successors, including:
\begin{table}[H]
  \centering
  \caption{\bert{} and contemporary models: a comparison of key properties.
           All models use the Transformer encoder architecture; differences lie in
           pre-training objectives, corpus size, and downstream performance.}
  \label{tab:model_comparison}
  \renewcommand{\arraystretch}{1.3}
  \begin{tabular}{lcccc}
    \toprule
    \textbf{Model} & \textbf{Type} & \textbf{Bidirectional} & \textbf{GLUE Score} & \textbf{Year} \\
    \midrule
    ELMo & Feature-based & Shallow & 76.2 & 2018 \\
    \bert{}-Base & Pre-train/Fine-tune & Yes (MLM) & 78.3 & 2018 \\
    GPT & Pre-train/Fine-tune & No (causal) & 72.3 & 2018 \\
    RoBERTa & Pre-train/Fine-tune & Yes (dynamic MLM) & 80.5 & 2019 \\
    ALBERT & Pre-train/Fine-tune & Yes (param. sharing) & 75.3 & 2019 \\
    \bottomrule
  \end{tabular}
\end{table}
\begin{itemize}
  \item \textbf{RoBERTa}~\cite{liu2019roberta}: removes NSP, trains with dynamic
        masking and larger batches, and demonstrates that pre-training for longer
        substantially improves performance.
  \item \textbf{ALBERT}~\cite{lan2019albert}: introduces parameter sharing across
        layers and factorised embedding parameterisation to reduce model size.
  \item \textbf{DistilBERT}~\cite{sanh2019distilbert}: uses knowledge distillation
        to produce a 40\%-smaller, 60\%-faster model retaining 97\% of \bert{}'s
        GLUE score.
  \item \textbf{Domain-specific BERTs}: BioBERT, SciBERT, ClinicalBERT, and others
        apply the same pre-training procedure to specialised corpora, demonstrating
        the generality of the framework.
\end{itemize}

\subsection{Conclusion}
\label{subsec:conclusion}

\bert{} marks a watershed moment in the evolution of NLP.  Its introduction of
deep bidirectional pre-training via masked language modelling demonstrated
conclusively that a single general-purpose encoder, pre-trained at scale, can
outperform highly engineered task-specific architectures across the full spectrum
of language understanding tasks.  The architectural and training insights embedded
in \bert{} continue to underpin state-of-the-art large language models, and its
influence on the field remains profound.
